\section{Day 2: \texorpdfstring{$\pi$}{pi}-\texorpdfstring{$\lambda$}{lambda} systems and Dynkin's theorem (Sep.\ 14, 2026)}
{\color{crimson} Office hours will be held by Nitit (the TA) on Mondays, from 5 to 6pm.}

Today, we will talk about Dynkin's $\pi$-$\lambda$ theorem (cf.\ p.5 of Panchenko's notes).
\begin{definition} \label{defn:pi_lambda_system}
    A $\pi$-system is a collection of sets $\SP$ which is closed under finite intersection, i.e.,
    \[ A, B \in \SP \implies A \cap B \in \SP. \]
    A $\lambda$-system is a nonempty collection of sets $\mathcal L$ closed under complements and countable disjoint unions, i.e.,
    \[ A \in \mathcal L \implies A^c \in \mathcal L, \]
    and if $(A_i : i \in \NN)$ is a countable collection of disjoint sets in $\mathcal L$, then their union $\bigcup_{i \in \NN} A_i \in \mathcal L$.
\end{definition}

\begin{theorem}[Dynkin] \label{thm:dynkin}
    If $\SP$ is a $\pi$-system and $\mathcal L$ is a $\lambda$-system such that $\SP \subset \mathcal L$, then $\sigma(\SP) \subset \mathcal L$.
\end{theorem}
Recall that we denote $\sigma(\SP)$ to be the $\sigma$-algebra generated by $\SP$, i.e., the smallest $\sigma$-algebra containing $\SP$ (obtained by taking the intersection over all such $\sigma$-algebras).

As an example of how one might use \ref{thm:dynkin}, let $\PP$ and $\QQ$ be probability measures on a measurable space $(\Omega, \mathcal S)$. Suppose, for some $\pi$-system $\SP$, that we have $\PP(A) = \QQ(A)$ for all $A \in \SP$. Then $\PP(A) = \QQ(A)$ for all $A \in \sigma(\SP)$.

\begin{proof}
    Let $\SA = \{A \in \mathcal S : \PP(A) = \QQ(A)\}$. $\SA$ is a $\lambda$-system, since
    \begin{align*}
        \PP(A) = \QQ(A) &\implies \PP(A^c) = 1 - \PP(A) \\
        &\PP(A^c) = 1 - \QQ(A) = \QQ(A^c),
    \end{align*}
    whence it is closed under complements; in particular, it is also closed under countable disjoint unions by countable additivity of $\PP$ and $\QQ$. Since $\SP \subset \SA$, it follows that \ref{thm:dynkin} would yield $\sigma(\SP) \subset \SA$ as desired.
\end{proof}

We now prove Dynkin's $\pi$-$\lambda$ theorem.
\begin{proof}[Proof (of \ref{thm:dynkin})]
    First, we want to show that ``$\pi + \lambda = \sigma$'', i.e., if $\SA$ is both a $\pi$-system and a $\lambda$-system, then it is a $\sigma$-algebra. It suffices to check that $\SA$ is closed under countable unions: if $(A_i : i \in \NN)$ is in $\SA$, then we may write
    \[ \bigcup_{i=1}^\infty A_i = \bigcup_{i=1}^\infty A_i \setminus \left(\bigcup_{j=1}^{i-1} A_j\right) = \bigcup_{i=1}^\infty \left(A_i \cap \bigcap_{j=1}^{i-1} A_j^c\right) \]
    to realize the union of $(A_i : i \in \NN)$ as a countable disjoint union, whence we are done per definition of $\lambda$-systems.

    Thus, it is enough to show that $\lambda(\SP)$ (which we will hereby denote as the $\lambda$-system generated by $\SP$) is a $\pi$-system, since from there, we can finish by arguing:
    \begin{itemize}
        \item $\lambda(\SP)$ exists, since being a $\lambda$-system is closed under arbitrary intersections.
        \item Then $\SP \subset \lambda(\SP) \subset \mathcal L$.
        \item From earlier, it would follow that $\lambda(\SP) = \sigma(\SP)$.
    \end{itemize}

    For a fixed set $A$, let $\SG_A = \{B \in \Omega : B \cap A \in \lambda(\SP)\}$. We will show that if $A \in \lambda(\SP)$, then $\lambda(\SP) \subset \SG_A$, in three steps as below:
    \begin{enumerate}[(i)]
        \item If $A \in \lambda(\SP)$, then $\SG_A$ is a $\lambda$-system. Let $(B_i : i \in \NN)$ be disjoint sets in $\SG_A$. Then
        \[ \left(\bigcup_{i=1}^\infty B_i\right) \cap A = \bigcup_{i=1}^\infty (B_i \cap A), \]
        for which each of $B_i \cap A$ is in $\lambda(\SP)$, since $B_i \in \SG_A$. For complements, let $B \in \SG_A$. Then
        \[ B^c \cap A = \left((B \cap A) \cup A^c\right)^c \in \lambda(\SP), \]
        so $B^c \in \SG_A$ as well. Thus, $\SG_A$ is a $\lambda$-system.
        \item Now, if $A \in \SP$, then $\lambda(\SP) \in \SG_A$. If we can show that $\SP \subset \SG_A$, then $\lambda(\SP) \subset \SG_A$, since $\SG_A$ is a $\lambda$-system by (i). Why is $\SP \subset \SG_A$? If $B \in \SP$, then $B \cap A \in \SP$, and so $B \cap A \in \lambda(\SP)$ implies $B \in \SG_A$.
        \item If $B \in \lambda(\SP)$, then $\lambda(\SP) \subset \SG_B$. By (ii), $\SP \subset \SG_B$, and by (i), $\SG_B$ is a $\lambda$-system. Since $\lambda(\SP)$ is the smallest $\lambda$-system (given by the intersection of all $\lambda$-systems containing $\SP$), we have that $\lambda(\SP) \subset \SG_B$. Note that this concludes the property that $\lambda(\SP)$ is a $\pi$-system, per definition of $\SG_B$. \qedhere
    \end{enumerate}
\end{proof}
For a more concise proof, see pg.\ 6 in Panchenko.